\documentclass[9pt, aspectratio=169]{beamer}
% \documentclass[10pt]{beamer}
\usepackage[utf8]{inputenc}
\usepackage[T1]{fontenc}
\usepackage[english]{babel}
\usetheme{Frankfurt}

\usepackage[backend=biber, style=authoryear]{biblatex}
\addbibresource{local_references.bib}

%\usepackage{lmodern}
\usepackage{amsfonts,amssymb,amsmath}
\usepackage[english]{babel}
\usetheme{Frankfurt}

\usepackage{csquotes}
\usepackage{setspace}

\usepackage{colortbl}
\usepackage{tabularx}
\renewcommand\tabularxcolumn[1]{m{#1}}

% --- Tickz
\usepackage{physics}
\usepackage{amsmath}
\usepackage{tikz}
\usepackage{mathdots}
\usepackage{yhmath}
\usepackage{cancel}
\usepackage{color}
\usepackage{siunitx}
\usepackage{array}
\usepackage{multirow}
\usepackage{amssymb}
\usepackage{gensymb}
\usepackage{tabularx}
\usepackage{extarrows}
\usepackage{booktabs}
\usetikzlibrary{fadings}
\usetikzlibrary{patterns}
\usetikzlibrary{shadows.blur}
\usetikzlibrary{shapes}

% ---------

\usepackage{booktabs}
\usepackage{setspace}
\usepackage{amssymb}
\usepackage{adjustbox}
\usepackage{pifont}
\usepackage[inkscapeformat=png]{svg}
\usepackage{graphicx}
\usepackage{times}
\setbeamertemplate{caption}[numbered]
% % \setbeamertemplate{bibliography item}{[\theenumiv]}

\setbeamerfont{bibliography item}{size=\tiny}
\setbeamerfont{bibliography entry author}{size=\tiny}
\setbeamerfont{bibliography entry title}{size=\tiny}
\setbeamerfont{bibliography entry location}{size=\tiny}
\setbeamerfont{bibliography entry note}{size=\tiny}

\setbeamerfont{frametitle}{size=\large}

\usepackage{caption}
\usepackage{float}
\usepackage{xcolor}
\usepackage{listings}
\usepackage{animate}

\definecolor{codegreen}{rgb}{0,0.6,0}
\definecolor{codegray}{rgb}{0.5,0.5,0.5}
\definecolor{codepurple}{rgb}{0.58,0,0.82}
\definecolor{backcolour}{rgb}{0.95,0.95,0.92}
 
\lstdefinestyle{mystyle}{
    backgroundcolor=\color{backcolour},   
    commentstyle=\color{codegreen},
    keywordstyle=\color{magenta},
    numberstyle=\tiny\color{codegray},
    stringstyle=\color{codepurple},
    basicstyle=\footnotesize,
    breakatwhitespace=false,         
    breaklines=true,                 
    captionpos=b,                    
    keepspaces=true,                 
    numbers=left,                    
    numbersep=5pt,                  
    showspaces=false,                
    showstringspaces=false,
    showtabs=false,                  
    tabsize=2
}
 
\lstset{style=mystyle}


\usepackage{ragged2e}
\setbeamercolor{section in foot}{fg=white,bg=darkorange}
\setbeamercolor{subsection in foot}{fg=white,bg=darkorange}
\setbeamercolor{frametitle}{fg=white, bg=darkorange}
\setbeamercolor{title}{fg=white, bg=darkorange}
\setbeamercolor{frame}{bg=darkorange}
\setbeamercolor{block title}{bg=darkorange,fg=white}

\setbeamercolor{item}{fg=darkorange}

\definecolor{darkorange}{rgb}{0.81, 0.52, 0.05}
\definecolor{darkorange2}{rgb}{1, 0.64, 0.2}
\definecolor{honeydew}{rgb}{1, 0.85, 0.45}


\newenvironment{variableblock}[3]{%
  \setbeamercolor{block body}{#2}
  \setbeamercolor{block title}{#3}
  \begin{block}{#1}}{\end{block}}

\newenvironment{prosblock}[1]{%
  % \setbeamercolor{block body}{bg=blue,fg=white}
  \setbeamercolor{block title}{bg=blue,fg=white}
  \begin{block}{#1}}{\end{block}}

\newenvironment{consblock}[1]{%
  % \setbeamercolor{block body}{bg=red,fg=white}
  \setbeamercolor{block title}{bg=red,fg=white}
  \begin{block}{#1}}{\end{block}}

\newcommand{\cmark}{\ding{51}}%
\newcommand{\xmark}{\ding{55}}%

\renewcommand{\arraystretch}{1.5}

% Please add the following required packages to your document preamble:
\usepackage{booktabs}
\usepackage{multirow}
\usepackage{colortbl}
% Beamer presentation requires \usepackage{colortbl} instead of \usepackage[table,xcdraw]{xcolor}

\usepackage{tabularray}\UseTblrLibrary{varwidth}
\usepackage{xcolor}
\def\BibTeX{{\rm B\kern-.05em{\sc i\kern-.025em b}\kern-.08em
    T\kern-.1667em\lower.7ex\hbox{E}\kern-.125emX}}
% \usepackage{cite}
\usepackage{amsmath}
\newcommand{\probP}{\text{I\kern-0.15em P}}
\usepackage{etoolbox}
\patchcmd{\thebibliography}{\section*{\refname}}{}{}{}

\setlength\tabcolsep{0.5pt}

\renewcommand{\arraystretch}{0.9}
\setlength{\tabcolsep}{2pt}

\usepackage{pgffor}

\setbeamerfont{bibliography item}{size=\tiny}
\setbeamerfont{bibliography entry author}{size=\tiny}
\setbeamerfont{bibliography entry title}{size=\tiny}
\setbeamerfont{bibliography entry location}{size=\tiny}
\setbeamerfont{bibliography entry note}{size=\tiny}

\setbeamerfont{bibliography entry author}{shape=\upshape,series=\mdseries,size=\footnotesize}
\setbeamerfont{bibliography entry title}{shape=\slshape,series=\mdseries,size=\footnotesize}
\setbeamerfont{bibliography entry journal}{shape=\upshape,series=\mdseries,size=\footnotesize}
\setbeamerfont{bibliography entry note}{shape=\upshape,series=\mdseries,size=\footnotesize}

\renewcommand*{\bibfont}{\scriptsize}

\newenvironment<>{varblock}[2][.9\textwidth]{%
  \setlength{\textwidth}{#1}
  \begin{actionenv}#3%
    \def\insertblocktitle{#2}%
    \par%
    \usebeamertemplate{block begin}}
  {\par%
    \usebeamertemplate{block end}%
  \end{actionenv}}

% \setbeamertemplate{footline}[frame number]

\setbeamertemplate{footline}{
  \leavevmode%
  \hfill
  \usebeamercolor[fg]{page number in head/foot}%
  \scriptsize\insertframenumber/\inserttotalframenumber%
  \hspace{1em}
}

\newcommand{\BaselineAnim}[2][0.47\linewidth]{%
  \animategraphics[autoplay,loop,width=#1]{8}{figures/baselines/#2/frame}{0}{199}%
}

\newcommand{\ComparisonAnimations}[2]{%
  \begin{columns}[T,totalwidth=\textwidth]
    \begin{column}{0.49\textwidth}
      \centering
      {\footnotesize\textbf{Reference: ORBITAL-MOISE+MARL}}\\[-0.05cm]
      \BaselineAnim[0.92\linewidth]{ORBITAL-MOISE+MARL}
    \end{column}
    \begin{column}{0.49\textwidth}
      \centering
      {\footnotesize\textbf{#2}}\\[-0.05cm]
      \BaselineAnim[0.92\linewidth]{#1}
    \end{column}
  \end{columns}
}

\newcommand{\CompactItem}[1]{\item \scriptsize #1}

\newcommand{\ExplainerImage}[2][0.92\linewidth]{%
  \includegraphics[width=#1]{figures/explainer/#2}%
}


\begin{document}

\author{
  \textbf{Julien Soulé}$^{1}$}

\title{\textbf{Presentation of the ORBITAL project and results}}

\subtitle{Presentation}

\institute{
  {\phantom{x}$^{1}$ \footnotesize \textit{University of Luxembourg, SnT, Luxembourg, Luxembourg}} \\ \phantom{x}}


\date{\textit{\footnotesize \today}}

%\subject{}
\setbeamercovered{transparent}
\setbeamertemplate{navigation symbols}{}
\begin{frame}[plain]
  \maketitle\vspace{-0.8cm}

  \ \\

  \begin{minipage}{\textwidth}
    \centering
    \textit{\href{mailto:julien.soule@uni.lu}{julien.soule@uni.lu}}
  \end{minipage}


\end{frame}

\begin{frame}{First: what is the mission loop?}
  \scriptsize
  \begin{columns}[T,totalwidth=\textwidth]
    \begin{column}{0.48\textwidth}
      \centering
      \ExplainerImage[0.84\linewidth]{environment_overview.png}
    \end{column}
    \begin{column}{0.50\textwidth}
      \begin{block}{What agents are trying to do}
        \begin{enumerate}
          \item \textbf{Discover/acquire:} discover known-local tasks, then observe yellow/or tasks near their orbit.
          \item \textbf{Carry:} store acquired data in their local buffer.
          \item \textbf{Deliver:} use REL\_GRN for direct ground delivery or REL\_SAT for satellite proxy routing.
          \item \textbf{Stabilize:} preserve energy, health, connectivity, cyber state, and debris safety.
        \end{enumerate}
      \end{block}
      \begin{block}{Most important idea}
        Observing a task is not enough. A satellite can collect data and still fail the mission if that data remains in its buffer and never reaches ground.
      \end{block}
    \end{column}
  \end{columns}
\end{frame}

\begin{frame}[shrink=8]{Visual grammar: what does each symbol mean?}
  \scriptsize
  \begin{columns}[T,totalwidth=\textwidth]
    \begin{column}{0.31\textwidth}
      \centering
      \ExplainerImage[0.70\linewidth]{satellite_marker.png}
      \begin{block}{Satellite marker}
        Number = satellite id. Cyan top bar = buffered data. Green lower bar = energy. Red lower bar = health. Text labels show role and last action.
      \end{block}
    \end{column}
    \begin{column}{0.31\textwidth}
      \centering
      \ExplainerImage[0.70\linewidth]{observation_task.png}
      \begin{block}{Observation task}
        Yellow/or dots are known active tasks. Brighter/larger means higher priority. A successful OBS action removes the task and fills a buffer.
      \end{block}
    \end{column}
    \begin{column}{0.31\textwidth}
      \centering
      \ExplainerImage[0.58\linewidth]{mission_panel.png}
      \begin{block}{Mission panel}
        Watch alive satellites, isolated satellites, delivered value, and reward components to understand the global effect.
      \end{block}
    \end{column}
  \end{columns}
\end{frame}

\begin{frame}[shrink=8]{Connectivity: satellite links vs ground routes}
  \scriptsize
  \begin{columns}[T,totalwidth=\textwidth]
    \begin{column}{0.48\textwidth}
      \centering
      \ExplainerImage[0.82\linewidth]{communication_links.png}
      \begin{block}{Blue communication links}
        Blue lines mean satellites can communicate at this step. They are not deliveries by themselves; they are possible relay paths.
      \end{block}
    \end{column}
    \begin{column}{0.48\textwidth}
      \centering
      \ExplainerImage[0.40\linewidth]{ground_downlink.png}
      \begin{block}{Green ground-route cue}
        A green square is a ground station. A green ground-route line means buffered data currently has a feasible path to ground. REL\_GRN delivers through direct ground contact; REL\_SAT moves data and task knowledge through satellite links.
      \end{block}
    \end{column}
  \end{columns}
\end{frame}

\begin{frame}{Labels: role is not exactly action}
  \scriptsize
  \begin{columns}[T,totalwidth=\textwidth]
    \begin{column}{0.45\textwidth}
      \centering
      \ExplainerImage[0.78\linewidth]{satellite_marker.png}
    \end{column}
    \begin{column}{0.53\textwidth}
      \begin{block}{Role labels: intended responsibility}
        \begin{tabularx}{\linewidth}{>{\bfseries}l X}
          OBS  & Observer role: prioritize acquiring task data.                \\
          REL  & Relay role: prioritize delivery and communication utility.    \\
          SAFE & Safety role: preserve energy, cyber health, or debris safety. \\
          FREE & No organizational role constraint.                            \\
          SOFT & Reward-shaped guidance, but no hard role enforcement.         \\
        \end{tabularx}
      \end{block}
      \begin{block}{Action labels: what was just executed}
        \begin{tabularx}{\linewidth}{>{\bfseries}l X}
          OBS  & Observe a known local task. \quad REL\_GRN: exchange with ground. \quad REL\_SAT: exchange with a satellite. \\
          PWR  & LowPower / recharge-oriented safe mode.                                         \\
          SCAN & CyberScan. \quad UP/DN: orbit maneuver.                                         \\
          IDLE & Do nothing mission-active this step; keep state/buffer, spend only idle energy. \\
        \end{tabularx}
      \end{block}
    \end{column}
  \end{columns}
\end{frame}

\begin{frame}{Action transition: Observe}
  \scriptsize
  \begin{columns}[T,totalwidth=\textwidth]
    \begin{column}{0.33\textwidth}
      \centering
      \ExplainerImage[0.70\linewidth]{observation_task.png}
      \begin{block}{Before}
        A nearby active task is visible as a warm dot. The satellite must be close enough to sense it.
      \end{block}
    \end{column}
    \begin{column}{0.30\textwidth}
      \begin{block}{Action}
        \textbf{OBS / Observe}: the satellite spends energy to service one local task.
      \end{block}
      \begin{block}{Immediate effect}
        The task becomes inactive and the satellite's cyan buffer increases.
      \end{block}
    \end{column}
    \begin{column}{0.33\textwidth}
      \centering
      \ExplainerImage[0.78\linewidth]{satellite_marker.png}
      \begin{block}{After}
        The data is only \textit{stored}. It is not mission value yet. Delivery still requires REL\_GRN at ground contact, possibly after REL\_SAT proxy routing.
      \end{block}
    \end{column}
  \end{columns}
\end{frame}

\begin{frame}[shrink=5]{Action transition: Relay}
  \scriptsize
  \begin{columns}[T,totalwidth=\textwidth]
    \begin{column}{0.34\textwidth}
      \centering
      \ExplainerImage[0.70\linewidth]{satellite_marker.png}
      \begin{block}{Before}
        A satellite has buffered data: cyan bar above the marker.
      \end{block}
    \end{column}
    \begin{column}{0.30\textwidth}
      \begin{block}{Action}
        \textbf{REL\_GRN / Relay ground}: deliver data and exchange task knowledge with a ground station. \textbf{REL\_SAT / Relay satellite}: move data and task knowledge to a better-positioned neighbor.
      \end{block}
      \begin{block}{Requirement}
        REL\_GRN requires direct ground contact. REL\_SAT requires a live communication link to another satellite.
      \end{block}
    \end{column}
    \begin{column}{0.32\textwidth}
      \centering
      \ExplainerImage[0.48\linewidth]{ground_downlink.png}
      \begin{block}{After}
        If ground contact is feasible, REL\_GRN decreases the buffer and increases delivered total. REL\_SAT can decrease one buffer while increasing another closer to ground.
      \end{block}
    \end{column}
  \end{columns}
\end{frame}

\begin{frame}{Action transitions: safety and waiting}
  \scriptsize
  \begin{columns}[T,totalwidth=\textwidth]
    \begin{column}{0.38\textwidth}
      \centering
      \ExplainerImage[0.66\linewidth]{isolation_task_pressure.png}
      \begin{block}{Stress cues}
        Yellow dashed ring = isolated; red ring = malware compromise; purple/magenta halos = debris. These are reasons to stop pure acquisition.
      \end{block}
    \end{column}
    \begin{column}{0.58\textwidth}
      \begin{block}{Before $\rightarrow$ action $\rightarrow$ after}
        \begin{tabularx}{\linewidth}{>{\bfseries}l X}
          PWR   & LowPower: spend little energy and recover when the satellite is sunlit; buffer remains stored.                              \\
          SCAN  & CyberScan: reduce/clear malware compromise; look for red compromise rings disappearing over time.                           \\
          UP/DN & OrbitUp/OrbitDown: maneuver to reduce local debris conjunction risk; radius changes between shells.                                  \\
          IDLE  & No active mission operation. The satellite keeps its buffer and position evolves, but it does not observe, relay, scan, or maneuver. \\
        \end{tabularx}
      \end{block}
      \begin{block}{Why IDLE after Observe?}
        After OBS, the satellite may wait because it has no task nearby, no delivery path yet, a role constraint says to hold, or another action would waste energy. IDLE does not erase the buffer; it simply delays delivery.
      \end{block}
    \end{column}
  \end{columns}
\end{frame}

\begin{frame}{How to read a policy live}
  \scriptsize
  \begin{columns}[T,totalwidth=\textwidth]
    \begin{column}{0.48\textwidth}
      \begin{block}{Ask these questions while watching}
        \begin{enumerate}
          \item Are satellites observing tasks or ignoring them?
          \item Are cyan buffers growing without delivery?
          \item Do REL\_GRN actions happen at ground contact and REL\_SAT actions route data through neighbors?
          \item Do SAFE/PWR/SCAN/UP/DN actions appear when stress appears?
          \item Are isolated satellites still observing myopically?
        \end{enumerate}
      \end{block}
    \end{column}
    \begin{column}{0.48\textwidth}
      \begin{block}{Good behavior in one sentence}
        The fleet should discover tasks, acquire data when tasks are valuable, route it to ground, and switch to safety actions before energy, health, cyber, debris, or isolation collapses the mission.
      \end{block}
      \begin{block}{What the comparisons show}
        The baselines differ mostly in when they switch between acquisition, delivery, and stabilization.
      \end{block}
    \end{column}
  \end{columns}
\end{frame}

\begin{frame}{ORBITAL in one operational view}
  \begin{columns}[T,totalwidth=\textwidth]
    \begin{column}{0.60\textwidth}
      \centering
      \BaselineAnim[0.94\linewidth]{ORBITAL-MOISE+MARL}
    \end{column}
    \begin{column}{0.38\textwidth}
      \scriptsize
      \begin{block}{What to read on the animation}
        \begin{itemize}
          \item Earth is centered; the three rings are LEO shells.
          \item Numbered satellites carry role/action labels plus buffer, energy, and health bars.
          \item Yellow/or dots are known observation tasks; green squares are ground stations.
          \item Blue links show inter-satellite communication; green ground-route lines show delivery opportunities.
          \item Purple/magenta halos indicate debris risk; red rings indicate malware compromise; yellow dashed rings indicate isolation.
        \end{itemize}
      \end{block}
      \begin{block}{Benchmark goal}
        Maximize useful Earth-observation data delivered to ground while preserving fleet health, connectivity, and resilience.
      \end{block}
    \end{column}
  \end{columns}
\end{frame}

\begin{frame}{Acquire--deliver--stabilize doctrine}
  \begin{columns}[T,totalwidth=\textwidth]
    \begin{column}{0.48\textwidth}
      \centering
      \BaselineAnim[0.98\linewidth]{ORBITAL-MOISE+MARL}
    \end{column}
    \begin{column}{0.50\textwidth}
      \scriptsize
      \begin{block}{Why the task is non-trivial}
        \begin{itemize}
          \item \textbf{Discover/acquire:} satellites learn dynamic high-value targets, observe known reachable tasks, and store data locally.
          \item \textbf{Deliver:} data only becomes mission value when REL\_GRN reaches a ground station, often after REL\_SAT proxy routing.
          \item \textbf{Stabilize:} energy, health, malware compromise, isolation, and debris risk can collapse the mission if ignored.
        \end{itemize}
      \end{block}
      \begin{block}{Core design feature}
        Observation and delivery are decoupled. A policy can look locally productive while failing globally if buffers never reach ground.
      \end{block}
    \end{column}
  \end{columns}
\end{frame}

\begin{frame}{Baseline ladder: from rules to organization}
  \scriptsize
  \begin{table}
    \centering
    \begin{tabularx}{0.96\textwidth}{>{\bfseries}l X X}
      \toprule
      Baseline           & Main visual signature                             & Main limitation                                                 \\
      \midrule
      RB-Rule            & OBS-heavy myopic acquisition                       & Late REL\_GRN, little proxy routing, weak stress adaptation      \\
      RB-Relay-heavy     & REL\_GRN whenever possible, frequent REL\_SAT      & Conservative, under-observes tasks                              \\
      PB-DCOP-lite       & Periodic OBS/REL/GRN/SAFE blocks                   & Role plans lag when links/tasks/cyber states shift              \\
      LB-Unconstrained   & FREE noisy action mix                              & Active but unsafe and weakly explainable                        \\
      LB-RewardOnly      & SOFT score-driven blend of all actions             & Incentives guide but do not enforce structure                   \\
      LB-ActionOnly      & Rigid OBS/SAT/GRN/SAFE masks                       & Disciplined but can idle or miss mission value                  \\
      ORBITAL-MOISE+MARL & Phase-aware discover/acquire/deliver/resilience    & Structured reference controller                                 \\
      \bottomrule
    \end{tabularx}
  \end{table}
  \vspace{0.15cm}
  \begin{center}
    \small Reference comparison strategy: every following slide keeps \textbf{ORBITAL-MOISE+MARL on the left} and the comparison baseline on the right.
  \end{center}
\end{frame}

\begin{frame}{Reference vs RB-Rule}
  \ComparisonAnimations{RB-Rule priority-first}{RB-Rule priority-first}
  \vspace{-0.1cm}
  \begin{block}{Synthetic comparison}
    \begin{itemize}
      \CompactItem{MOISE+MARL shifts between discovery, acquisition, delivery, and resilience; RB-Rule mostly keeps pressing OBS.}
      \CompactItem{In RB-Rule, look for OBS persistence and cyan buffer growth before late REL\_GRN attempts.}
      \CompactItem{Key message: interpretability is high, but the policy is myopic and reacts late to delivery, energy, cyber, and debris pressure.}
    \end{itemize}
  \end{block}
\end{frame}

\begin{frame}{Reference vs RB-Relay-heavy}
  \ComparisonAnimations{RB-Relay-heavy}{RB-Relay-heavy}
  \vspace{-0.1cm}
  \begin{block}{Synthetic comparison}
    \begin{itemize}
      \CompactItem{MOISE+MARL relays when delivery is valuable, but still keeps observers active when known tasks are reachable.}
      \CompactItem{RB-Relay-heavy should show many GRN/REL labels and REL\_GRN/REL\_SAT actions, often before enough data is collected.}
      \CompactItem{Key message: immediate delivery discipline can starve acquisition and waste relay opportunities.}
    \end{itemize}
  \end{block}
\end{frame}

\begin{frame}{Reference vs PB-DCOP-lite}
  \ComparisonAnimations{PB-DCOP-lite}{PB-DCOP-lite}
  \vspace{-0.1cm}
  \begin{block}{Synthetic comparison}
    \begin{itemize}
      \CompactItem{MOISE+MARL adapts roles continuously through organizational guides; PB-DCOP-lite freezes OBS/REL/GRN/SAFE blocks for several steps.}
      \CompactItem{PB-DCOP-lite should look clean and blocky: ground relayers, satellite relayers, observers, and safety agents.}
      \CompactItem{Key message: planning improves structure, but assignments can lag behind fast link drops, cyber events, and debris drift.}
    \end{itemize}
  \end{block}
\end{frame}

\begin{frame}{Reference vs LB-Unconstrained}
  \ComparisonAnimations{LB-Unconstrained}{LB-Unconstrained}
  \vspace{-0.1cm}
  \begin{block}{Synthetic comparison}
    \begin{itemize}
      \CompactItem{MOISE+MARL constrains behavior through roles and safety overrides; LB-Unconstrained keeps full action freedom.}
      \CompactItem{On the right, FREE labels indicate agents can observe, relay, idle, maneuver, power, or scan without organizational discipline.}
      \CompactItem{Key message: unconstrained learning can be active and adaptive, but it is brittle, less auditable, and prone to unsafe persistence.}
    \end{itemize}
  \end{block}
\end{frame}

\begin{frame}{Reference vs LB-RewardOnly}
  \ComparisonAnimations{LB-RewardOnly}{LB-RewardOnly}
  \vspace{-0.1cm}
  \begin{block}{Synthetic comparison}
    \begin{itemize}
      \CompactItem{MOISE+MARL combines reward guidance with role-action control; RewardOnly samples from soft scores for OBS, REL\_GRN, REL\_SAT, and safety.}
      \CompactItem{The SOFT label means the policy is encouraged to balance task, delivery, and safety, but actions remain unconstrained.}
      \CompactItem{Key message: reward shaping improves alignment, yet cannot guarantee role consistency or safety compliance by itself.}
    \end{itemize}
  \end{block}
\end{frame}

\begin{frame}{Reference vs LB-ActionOnly}
  \ComparisonAnimations{LB-ActionOnly}{LB-ActionOnly}
  \vspace{-0.1cm}
  \begin{block}{Synthetic comparison}
    \begin{itemize}
      \CompactItem{MOISE+MARL uses flexible phase-aware roles; ActionOnly enforces fixed OBS/SAT/GRN/SAFE action masks without mission shaping.}
      \CompactItem{The right side should show rigid specialization, including idling when a fixed role cannot act.}
      \CompactItem{Key message: masks improve controllability and explainability, but can become rigid when mission value requires flexible role switching.}
    \end{itemize}
  \end{block}
\end{frame}

\begin{frame}{Overall interpretation}
  \scriptsize
  \begin{columns}[T,totalwidth=\textwidth]
    \begin{column}{0.46\textwidth}
      \begin{block}{What improves across the ladder}
        \begin{itemize}
          \item From local rules to structured role allocation.
          \item From raw action freedom to auditable role-action discipline.
          \item From immediate task servicing to long-horizon acquire--deliver--stabilize behavior.
        \end{itemize}
      \end{block}
    \end{column}
    \begin{column}{0.50\textwidth}
      \begin{block}{Why ORBITAL-MOISE+MARL is the reference}
        \begin{itemize}
          \item Better delivery continuity without abandoning acquisition.
          \item Safety reactions are explicit and visually explainable.
          \item Role specialization appears at trajectory level, not only as aggregate reward.
        \end{itemize}
      \end{block}
    \end{column}
  \end{columns}
  \vspace{0.4cm}
  \begin{center}
    \Large ORBITAL makes coordination failures visible, not only measurable.
  \end{center}
\end{frame}



\end{document}
